\documentclass[../main.tex]{subfiles}
\begin{document}

% Target length: <= 10 pages.

Chapter~\ref{chapter:Related_works} defined the system requirements: grounded answers, conflict-tolerant recommendations, graceful degradation under model failures, traceable advisory runs, and a responsive chat interface.
This chapter explains the technical choices made to satisfy those requirements.
Each section follows the same decision logic: the system need, the selected approach, the main alternative, and the reason for the selection.
Section~\ref{section:3.1} covers hosted \gls{LLM}s and structured output, Section~\ref{section:3.2} covers retrieval-augmented generation and vector search, Section~\ref{section:3.3} covers the fixed advisory graph, Section~\ref{section:3.4} covers document processing and \gls{OCR}, and Section~\ref{section:3.5} covers the web stack.

\section{Large language models and structured output}
\label{section:3.1}

The system needs language-model behavior in places where deterministic parsing is too brittle: routing a free-form student message, extracting profile updates from conversational text, resolving selected source conflicts, generating grounded answers, and transcribing scanned pages.
These tasks do not require training a new model; they require reliable calls to a capable hosted model with clear input and output contracts.

The selected approach is to use the Gemini family through a shared inference gateway~\cite{gemini,geminiteam2023}.
Gemini's multimodal interface is useful because the same gateway can handle both text tasks and page-image \gls{OCR}.
Each call site sends a role-specific prompt and, when downstream code consumes the result, requires strict \gls{JSON} rather than free prose.
The returned text is parsed into typed service models before it can affect routing, profile state, or answer generation.

The main alternative was to call a hosted model directly from every service or to run a local open-weight model.
Direct calls would duplicate retry, telemetry, and parsing logic across the codebase.
A local model would reduce provider dependence, but it would add deployment and hardware cost outside the project's scope.
The gateway-based hosted approach keeps the model boundary explicit: hard \gls{API} failures and malformed structured outputs are detected in one place, retries and fallback models are applied consistently, and deterministic fallbacks can be used by call sites when the model cannot return a usable result.

\section{Retrieval-augmented generation and vector search}
\label{section:3.2}

The knowledge-answering route needs current, source-grounded answers about tuition, programs, scholarships, regulations, and school-specific documents.
A bare \gls{LLM} cannot satisfy this need because the relevant admission material changes during the admission season and may never have appeared in the model's training data.
Retrieval-augmented generation addresses this by retrieving relevant passages from an external corpus and constraining generation to those passages~\cite{lewis2020rag}.

The selected approach stores document chunks as embeddings and retrieves them by semantic similarity.
Each chunk and each question is mapped to a dense vector; semantically related texts are close under cosine distance.
PostgreSQL with pgvector stores these vectors and supports nearest-neighbor search through the cosine-distance operator~\cite{postgresql,pgvector}.
The knowledge service embeds the question, searches the scoped chunk set, passes the retrieved chunks to the model, and attaches citations to the source documents used in the answer.

The main alternative was a dedicated vector database.
That choice would be reasonable for a much larger corpus or for specialized vector operations, but it would add another datastore to deploy and keep consistent with the relational admission data.
The thesis system already depends on PostgreSQL for sessions, traces, canonical records, cutoff records, and migrations.
Keeping vectors in the same database is sufficient for the measured corpus size recorded in \texttt{latex/FACTS.md} and lets advisory data, knowledge chunks, and traceable session state remain in one operational store.

\section{Agentic pipelines}
\label{section:3.3}

The advisory route needs a multi-stage decision process, not a single model answer.
The system must collect or update a profile, retrieve candidate programs, detect source conflicts, reason about eligibility, apply policy constraints, and explain the result.
The order of these stages should be predictable because recommendations are inspected by operators and may be re-run after a student correction.

The selected approach is a fixed LangGraph state graph~\cite{langgraph,langgraph_docs}.
Nodes read and write a shared typed state, and edges define the six-stage sequence: profile, retrieval, conflict, reasoning, policy, and explanation.
This structure gives every run the same stage names and the same control flow, which is why trace events can be persisted per stage and rendered in the operator panel.

The main alternative was an open-ended function-calling loop in which a model chooses tools until it decides that the task is complete.
That style is flexible, but the model controls the run-time sequence, making failures harder to reproduce and stage-level traces harder to compare.
For admission advising, the task shape is stable and the risk of an opaque tool sequence is higher than the value of improvisation.
The fixed graph was therefore chosen for determinism, traceability, and testable service boundaries.

\section{Document processing}
\label{section:3.4}

The ingestion layer must process official sources that arrive as \gls{HTML} pages, born-digital \gls{PDF}s, and scanned \gls{PDF}s.
Machine-readable pages can be parsed directly, but scanned proposal documents have no text layer and require \gls{OCR}.
The system therefore needs a document path that recovers usable text without making scanned files a separate project.

The selected approach is hybrid extraction.
For born-digital pages, the pipeline extracts the embedded text and tables using pdfplumber~\cite{pdfplumber}.
For scanned pages, it renders page images with PyMuPDF~\cite{pymupdf} and sends only those image pages through the multimodal inference gateway for transcription.
This keeps the cost proportional to the number of scanned pages and reuses the same hosted Gemini integration used by the other model-backed services.

The main alternative was a separate OCR engine.
That would reduce dependence on the hosted model for transcription, but it would add installation, language tuning, image preprocessing, and a second failure surface.
Vietnamese admission documents also mix tables, diacritics, scanned signatures, and layout artifacts, so the practical cost of tuning a separate OCR path is high.
Reusing the gateway is simpler for this project, while degenerate-output detection and retry protect the pipeline from the main observed model-OCR failure mode.

\section{Web stack}
\label{section:3.5}

The web layer needs to accept a student message immediately while advisory work continues in the background.
It also needs typed boundaries between routes, services, repositories, and inference calls so that malformed input or model output does not spread through the system.

The selected approach is FastAPI served by Uvicorn, with server-rendered Jinja2 templates and client-side polling for run completion~\cite{fastapi,uvicorn,jinja2}.
The message endpoint records the turn and returns quickly; the advisory run executes on a background thread pool and writes its result and trace to the database.
Pydantic models define the service contracts used by the chat layer, inference gateway, knowledge service, and profile state handling~\cite{pydantic}.

The main alternative was to introduce a separate task queue and worker service.
That would be appropriate for high-volume production traffic, but it would add a broker and another deployment unit.
For this thesis system, a small number of long, \gls{API}-bound runs can be handled by the existing process, while the database remains the durable handoff point between request handling, background execution, polling, and trace inspection.

This chapter connected each technology choice to a requirement and an alternative.
The system uses hosted Gemini calls behind a resilient gateway, \gls{RAG} over pgvector for grounded answers, a fixed LangGraph graph for deterministic advisory runs, hybrid extraction for mixed document formats, and an asynchronous FastAPI stack for non-blocking chat.
Chapter~\ref{chapter:Experiment} next describes how these choices are assembled into the implemented system.

\end{document}
